\section{Introduction and Background}
\label{sec:introduction}

\subsection{NegDRO and causal invariance}

Multi-environment data can distinguish a stable causal relationship from
environment-specific predictive associations.  Wang, Hu, B\"uhlmann, and
Guo~\cite{wang2026negdro} formulate this idea as Negative-Weighted
Distributionally Robust Optimization (NegDRO), a continuous minimax problem
whose effective environmental weights may be negative.  The negative weights
penalize disparities between environmental risks and thereby encourage risk
invariance.  Under additive interventions and sufficient environmental
heterogeneity, their population analysis relates stationarity of the NegDRO
objective to recovery of the causal outcome coefficient $\beta^*$.

The official algorithms in arXiv:2412.11850v3 optimize empirical risks from a
fixed data set.  In particular, official Algorithm~1 computes an exact
penalized maximizer in the environmental weight and then takes a primal
gradient step.  The present work addresses a different question: can one use
one fresh environmental observation per round, update the primal and dual
variables simultaneously, and obtain a direct population error bound?  The
answer given here is conditional on an explicit fresh-sampling law and moment,
geometry, and heterogeneity assumptions.

\subsection{The stochastic extension}

At each round, our method evaluates both stochastic oracles at the current
predictable state $(b_t,w_t)$.  It takes a projected descent step in $b$ and an
exponentiated-gradient ascent step in $w$.  The two oracles may be formed from
the same observation; their mutual independence is neither assumed nor used.
The analysis compares $w_t$ with a fixed simplex vector $w^0$ whose
heterogeneity matrix has eigenvalues bounded below by $\lambda>0$.  This
witness is an analytical comparator, not the algorithmic initialization.

The proof has a common structure in the penalized and unpenalized cases:
\[
  \begin{gathered}
    \text{corrected risk expansion}
    \Longrightarrow \text{fixed-witness curvature}
    \Longrightarrow \text{comparator identity}\\
    \Longrightarrow \text{entropy regret}
    \Longrightarrow \text{projected primal recursion}.
  \end{gathered}
\]
The unpenalized comparator identity transfers the dual regret with a factor of
two.  The penalized identity adds exactly the direction loss $2\mu$, which
becomes the squared-error bias $2\mu/\lambda$ after the curvature division.
Instantaneous comparator regret is not assumed to be nonnegative.

\subsection{Contributions and verified scope}

This note makes the following contributions.
\begin{itemize}[leftmargin=*]
  \item It states a corrected realized reduced form for additive interventions and
        derives the common population-risk expansion used by the online
        analysis.  The sign convention is forced by direct elimination from
        the structural equations; the discrepancy with the printed official
        Appendix~E.1 identity is documented in
        Section~\ref{subsec:corrected-reduced-form}.
  \item It analyzes a fresh-sample stochastic projected-SGD/EG recursion on a
        general nonempty closed convex bounded set.  The universal conditional
        sampling law yields the required selected moments, conditional oracle
        means, and oracle integrability.
  \item It gives exact finite-time bounds for the time-average of expected
        squared errors and for the expected squared error of the averaged
        iterate.  Equal steps $\eta_b=\eta_w=T^{-1/2}$ yield explicit
        inverse-square-root optimization terms.
  \item It provides a complete Lean-verified conditional proof of this chain.
        Completeness is relative to the stated model, sampling, moment,
        spectral, and feasible-set hypotheses; it is not an assumption-free
        claim or an external certification.
\end{itemize}

The result is neither a last-iterate nor a high-probability theorem.  A
separate square-root consequence concerns root mean square error and must not
be read as a bound on expected norm.  The project also does not construct an
infinite independent sample stream or a product probability space.  These
scope boundaries are summarized in Section~\ref{sec:discussion}.

We use arXiv:2412.11850v3, revised January~7, 2026, as the source version of
the original paper.  Statements attributed to that paper and statements
proved for the stochastic extension are distinguished throughout.
